% Z2 holonomy <-> Cech H2 obstruction and spectral splitting.

\paragraph{\texorpdfstring{$\ZZ_2$}{Z/2} 余圈 \texorpdfstring{$\leftrightarrow$}{<->} Čech-\texorpdfstring{$H^2$}{H2} 障碍：可见传输、最小双覆盖去奇异化与扭曲谱分裂}\label{par:xi-null-z2-cech-transfer-doublecover}
本段把“粘合扭结（定理 \ref{thm:xi-null-orthogonal-trichotomy-resource-nonsubstitutable} 第 (2) 类）”在一个最小可见通道 \(\mu_2\simeq\ZZ_2\) 上进一步刚性化：时间/协议图上的 \(\ZZ_2\) 余圈 holonomy 通过一个自然传输映射生成 \(\ZZ_2\)-值的 Čech-\(H^2\) 障碍；并且该 \(\ZZ_2\) 障碍在最小取向双覆盖上被拉回消去，对应扭曲通道在谱侧分裂为一对直和块。

\begin{theorem}[\texorpdfstring{$\ZZ_2$}{Z/2} 余圈--Čech 障碍的传输对应与 \texorpdfstring{$\mathsf{NULL}$}{NULL} 的不可规避性]\label{thm:xi-null-z2-cech-transfer-bockstein}
\leanverified{paper\_xi\_null\_z2\_cech\_transfer\_bockstein}
沿用定义 \ref{def:prefix-site-h2-gluing-obstruction} 的站点覆盖 \(\{b_i\succeq a\}\) 与粘合障碍类 \([\alpha]\in H^2(\{b_i\},A)\)。
再设时间/协议实现的有向图为 \(\Gamma\)，并在边上给定一条 \(\ZZ_2\)-余圈 \(\beta:E(\Gamma)\to\ZZ_2\)（其闭环符号
\(\chi_\beta(p)=(-1)^{\sum_{e\in p}\beta(e)}\)
见结论 \ref{con:xi-split-tunneling-zeta-detect-aflip}）。
假设系数阿贝尔群 \(A\) 含有唯一可见的二阶中心扭结 \(\mu_2=\{\pm1\}\subset A\)，并令自然投影
\[
\pi_2:A\twoheadrightarrow \mu_2\cong \ZZ_2.
\]
记 \(\ZZ_2\)-像
\[
[\alpha_2]:=\pi_2([\alpha])\in H^2(\{b_i\},\ZZ_2).
\]
则在“覆盖神经复形 \(1\)-骨架 \(\simeq \Gamma\)”的识别下，存在自然映射
\[
\mathrm{tr}:\ H^1(\Gamma;\ZZ_2)\longrightarrow H^2(\{b_i\},\ZZ_2),
\]
使得
\[
\mathrm{tr}([\beta])=[\alpha_2].
\]
因此
\(
[\beta]\neq 0 \iff [\alpha_2]\neq 0
\)；
并且若存在闭环 \(p\) 满足 \(\chi_\beta(p)=-1\)，则必有 \([\alpha_2]\neq 0\)，从而 \([\alpha]\neq 0\)，并由命题 \ref{prop:null-as-h2-obstruction} 推出 \(\mathscr{F}(a)=\varnothing\)，即该地址处出现真实的 \(\mathsf{NULL}\) 空纤维。
\end{theorem}

\begin{proof}
按定义 \ref{def:prefix-site-h2-gluing-obstruction}，对覆盖 \(\{b_i\}\) 上任意一族局部对象 \((s_i)\) 的二重重叠差异可由转移数据 \(g_{ij}\in A\) 记录，并诱导三重交上的 Čech \(2\)-余圈
\(
\alpha_{ijk}=g_{ij}+g_{jk}+g_{ki}
\)。
将 \(\pi_2\) 作用于转移数据并记 \(\bar g_{ij}:=\pi_2(g_{ij})\in\ZZ_2\)，则
\[
\bar g_{ij}+\bar g_{jk}+\bar g_{ki}=\alpha_{2,ijk}\qquad(\mathrm{mod}\ 2),
\]
即 \(\alpha_2=\delta(\bar g)\)。
另一方面，\(\ZZ_2\)-余圈 \([\beta]\in H^1(\Gamma;\ZZ_2)\) 可由边上的 \(\ZZ_2\) 标号刻画，而在神经复形 \(1\)-骨架与 \(\Gamma\) 的识别下，\(\bar g\) 给出一条 \(1\)-上链代表；其 Čech 余边界 \(\delta(\bar g)\) 正对应于所述传输映射 \(\mathrm{tr}\) 的像。
从而 \(\mathrm{tr}([\beta])=[\alpha_2]\)，并推出等价性。
最后，\(\chi_\beta(p)=-1\) 蕴含 \([\beta]\neq 0\)，故 \([\alpha_2]\neq 0\)；而 \([\alpha_2]\neq 0\Rightarrow [\alpha]\neq 0\) 由投影函子性立得，命题 \ref{prop:null-as-h2-obstruction} 给出 \(\mathscr{F}(a)=\varnothing\)。
证毕。
\end{proof}

\begin{corollary}[Čech 二障碍的单一来源：由取向单值化类完全决定]\label{cor:xi-null-z2-cech-obstruction-single-source-w1}
沿用定理 \ref{thm:xi-null-z2-cech-transfer-bockstein} 的记号，并把 \([\beta]\in H^1(\Gamma;\ZZ_2)\) 视为由符号单值化所诱导的取向类 \(w_1\)。
则有严格恒等式
\[
[\alpha_2]=\mathrm{tr}(w_1)\in H^2(\{b_i\},\ZZ_2).
\]
特别地，\([\alpha_2]=0\) 当且仅当 \(w_1=0\)（等价地：\([\beta]=0\)）。
\end{corollary}
\begin{proof}
由定理 \ref{thm:xi-null-z2-cech-transfer-bockstein} 得 \([\alpha_2]=\mathrm{tr}([\beta])\)。
将 \(w_1\) 记作 \([\beta]\) 的别名即得所示恒等式；最后一句为其直接推论。证毕。
\end{proof}

\begin{theorem}[取向双覆盖的最小去奇异化：拉回后 \texorpdfstring{$\ZZ_2$}{Z/2} 障碍消失]\label{thm:xi-null-z2-doublecover-min-desingularization}
\leanverified{paper\_xi\_null\_z2\_doublecover\_min\_desingularization}
沿用定理 \ref{thm:xi-null-z2-cech-transfer-bockstein} 的设定，并令
\(\pi:\Gamma_\beta\to\Gamma\)
为由同调类 \([\beta]\in H^1(\Gamma;\ZZ_2)\) 所对应的连通两层覆盖（等价于指数 \(2\) 子群 \(\ker(\pi_1(\Gamma)\twoheadrightarrow\ZZ_2)\) 的覆盖）。
则：
\begin{enumerate}
  \item \(\pi^*([\beta])=0\in H^1(\Gamma_\beta;\ZZ_2)\)；
  \item \(\pi^*([\alpha_2])=0\in H^2(\{\pi^{-1}(b_i)\},\ZZ_2)\)；
  \item 若进一步假设 \([\alpha]\) 的全部可核验扭结由 \(\mu_2\) 通道承载（即 \([\alpha]\) 在可见化意义下由 \([\alpha_2]\) 完全决定），则在该双覆盖上拉回相容层在可见通道中中性化，从而消除由 \(\ZZ_2\) 扭结触发的 \(\mathsf{NULL}\) 空纤维现象。
\end{enumerate}
此外，当 \([\beta]\neq 0\) 时，上述指数 \(2\) 覆盖在消除 \(\ZZ_2\) 障碍意义下为最小次数连通覆盖。
\end{theorem}

\begin{proof}
(1) 由覆盖的分类构造：\(\Gamma_\beta\) 正是使 \([\beta]\) 拉回为余边界的最小连通覆盖。
(2) 由定理 \ref{thm:xi-null-z2-cech-transfer-bockstein} 的自然性，
\(
\pi^*([\alpha_2])=\pi^*\mathrm{tr}([\beta])=\mathrm{tr}(\pi^*([\beta]))=0
\)。
(3) 在假设下，\(\ZZ_2\) 障碍消失即意味着可见通道中的二阶一致性扭结被严格消去，故对应的可见读出不再以 \(\mathsf{NULL}\) 方式失败。
最小性：若某连通覆盖 \(\Gamma'\to\Gamma\) 使 \(\beta\) 拉回为零，则其对应子群包含 \(\ker(\beta)\)；当 \([\beta]\neq 0\) 时 \([\pi_1(\Gamma):\ker(\beta)]=2\)，故次数至少为 \(2\)。
证毕。
\end{proof}

\begin{theorem}[双覆盖上的谱分裂：\texorpdfstring{$\beta$}{beta}-扭曲与非扭曲通道的直和分解]\label{thm:xi-null-z2-spectral-splitting-doublecover}
\leanverified{paper\_xi\_null\_z2\_spectral\_splitting\_doublecover}
设有限状态转移矩阵在 \(\ZZ_2\) 分裂标号下写作
\[
B(u)=B_0+uB_1,
\]
其中 \(B_0\) 收集 \(\beta=0\) 的跃迁、\(B_1\) 收集 \(\beta=1\) 的跃迁；并记 \(\beta\)-扭曲矩阵为
\[
B_\beta(u)=B_0-uB_1.
\]
令 \(\widetilde B(u)\) 为取向双覆盖 \(\Gamma_\beta\) 上的提升转移矩阵，则存在显式相似变换使
\[
\widetilde B(u)\ \sim\ B(u)\ \oplus\ B_\beta(u).
\]
从而对任意 \(z\) 有行列式分解
\[
\det(I-z\widetilde B(u))=\det(I-zB(u))\cdot \det(I-zB_\beta(u)).
\]
\end{theorem}

\begin{proof}
在 sheet 基下，覆盖提升的标准分块形式为
\[
\widetilde B(u)=
\begin{pmatrix}
B_0 & uB_1\\
uB_1 & B_0
\end{pmatrix}.
\]
取正交矩阵
\(
Q=\frac1{\sqrt2}\begin{pmatrix}I&I\\ I&-I\end{pmatrix}
\)，
直接计算得
\[
Q^{-1}\widetilde B(u)Q=
\begin{pmatrix}
B_0+uB_1 & 0\\
0 & B_0-uB_1
\end{pmatrix}
=B(u)\oplus B_\beta(u),
\]
行列式分解随之成立。
证毕。
\end{proof}

\begin{theorem}[语义共轭下的 \texorpdfstring{$\mathsf{NULL}$}{NULL}--临界模态刚性]\label{thm:xi-null-z2-conjugacy-rigidity}
设两套有限状态实现 \((\Gamma,\ell)\)、\((\Gamma',\ell')\) 表示同一协议，并存在滑动块共轭使分裂余圈在共轭下仅差余边界（结论 \ref{con:xi-split-tunneling-cohomology-conjugacy-invariance} 的设定）。
则：
\begin{enumerate}
  \item \([\beta]=0 \iff [\beta']=0\)，因而 \([\alpha_2]=0 \iff [\alpha_2'] =0\)；
  \item 因此由 \(\mu_2\) 通道触发的“\(\mathsf{NULL}\) 是否出现”在语义共轭下保持不变；
  \item 并且由定理 \ref{thm:xi-null-z2-spectral-splitting-doublecover} 所诱导的扭曲/未扭曲通道谱分裂（及其临界达界 \(|\mu|=\sqrt{\lambda}\) 上的相位驻留）亦为语义不变量。
\end{enumerate}
\end{theorem}

\begin{proof}
第 (1) 点即结论 \ref{con:xi-split-tunneling-cohomology-conjugacy-invariance} 所给出的闭环 \(\ZZ_2\) holonomy 不变性。
第 (2) 点由定理 \ref{thm:xi-null-z2-cech-transfer-bockstein} 的函子性与命题 \ref{prop:null-as-h2-obstruction}（经 \(\mu_2\) 可见化通道）推出。
第 (3) 点由共轭诱导的覆盖提升与相似变换保持谱数据立得。
证毕。
\end{proof}

\begin{corollary}[解析判据：\texorpdfstring{$\ZZ_2$}{Z/2} 障碍 \texorpdfstring{$\Leftrightarrow$}{<->} 临界圆上的 \texorpdfstring{$(-1)$}{-1} 相位极点]\label{cor:xi-null-z2-analytic-pole-criterion}
沿用定理 \ref{thm:xi-null-z2-spectral-splitting-doublecover} 的记号，并令 \(\lambda:=\rho(B(1))>1\)。
设扭曲通道的 \(\zeta\)-函数（或 Fredholm 行列式倒数）由
\[
\zeta_\beta(z,u):=\det(I-zB_\beta(u))^{-1}
\]
给出（与结论 \ref{con:xi-split-tunneling-z2-determinant-ratio} 的行列式口径一致）。
则以下命题等价：
\begin{enumerate}
  \item \([\alpha_2]\neq 0\)（从而 \([\alpha]\neq 0\) 并触发粘合型 \(\mathsf{NULL}\)，命题 \ref{prop:null-as-h2-obstruction}）；
  \item \([\beta]\neq 0\in H^1(\Gamma;\ZZ_2)\)（等价地存在闭环 \(p\) 使 \(\chi_\beta(p)=-1\)）；
  \item 归一化临界圆 \(|z|=\lambda^{-1/2}\) 上，\(\zeta_\beta(z,1)\) 在点 \(z=-\lambda^{-1/2}\) 处出现极点；等价地，\(B_\beta(1)\) 存在特征值 \(\mu=-\sqrt{\lambda}\)（相位为 \(-1\) 的达界模态）。
\end{enumerate}
\end{corollary}

\begin{proof}
(1)\(\Leftrightarrow\)(2) 由定理 \ref{thm:xi-null-z2-cech-transfer-bockstein}。
(2)\(\Leftrightarrow\)(3) 由 \(\zeta_\beta(z,1)=\det(I-zB_\beta(1))^{-1}\)：\(z=-\lambda^{-1/2}\) 为极点当且仅当 \(-\lambda^{-1/2}\) 为 \(\det(I-zB_\beta(1))\) 的零点，当且仅当 \(B_\beta(1)\) 具有特征值 \(\mu=(-\lambda^{-1/2})^{-1}=-\sqrt{\lambda}\)。
证毕。
\end{proof}

\begin{theorem}[有限阶可见扭结强迫 cyclotomic 相位与周期分解]\label{thm:xi-cech-torsion-forces-period}
\leanverified{paper\_xi\_cech\_torsion\_forces\_period}
设某一有限态实现的转移核在可见扭结通道上诱导出一族扭曲算子块 \(B_\chi\)（角色/表示通道），并设平凡通道的 Perron 根为 \(\lambda>1\)。
固定整数 \(r>1\) 并令 \(\zeta_r:=e^{2\pi i/r}\)。
若对某个通道 \(\chi\) 有归一化行列式
\[
\det\!\Bigl(I-\frac{t}{\lambda}B_\chi\Bigr)
\]
在单位圆上包含一个本原 \(r\) 次单位根零点（等价地，\(B_\chi\) 存在特征值 \(\lambda\zeta_r^{-1}\)），并且整体非负整数核 \(B\)（可取包含全部通道的正则 lift 核）为不可约，
则 \(B\) 的周期 \(p\) 满足 \(r\mid p\)，从而其状态空间规模满足
\[
|V(B)|\ \ge\ p\ \ge\ r.
\]
\end{theorem}
\begin{proof}
若 \(B\) 为不可约非负整数矩阵，则 Perron--Frobenius 周期理论断言：存在整数 \(p\ge 1\) 使得 \(B\) 的全部模长等于 \(\rho(B)\) 的特征值恰为
\[
\rho(B)\,e^{2\pi i k/p}\qquad (k=0,1,\dots,p-1),
\]
其中 \(p\) 等于底层有向图回路长度的最大公因子（周期）。
当某一通道块出现达界特征值 \(\lambda\zeta_r^{-1}\) 时，该相位必须属于上述根系，故 \(\zeta_r\) 必为 \(p\) 次单位根，等价于 \(r\mid p\)。
另一方面，周期分解将顶点集分成 \(p\) 个非空循环类，故 \(|V(B)|\ge p\)。
合并即得 \(|V(B)|\ge p\ge r\)。证毕。
\end{proof}

\endinput
